% C, C++

\subsection{Metrics Evaluation for C++}

\subsubsection{Introduction}

C++ represents a multi-paradigm language with a predominantly nominal and statically checked type system, enriched by templates, multiple inheritance, and low-level memory control. Its object model has evolved significantly from early descriptions in \cite{stroustrup1994design} to the formalized ISO standard \cite{cpp2017iso}. Prior studies emphasize that C++ combines abstraction mechanisms with explicit control over performance and representation, often at the cost of formal safety guarantees \cite{abrahams2004c++}. This section evaluates C++ using the defined metrics, grounding each numerical assignment in the language specification and established literature.

\subsubsection{Type System and Subtyping Model}

% TODO: Перепроверить сорсы

\textbf{1) Nominal–Structural Typing Score $S_{nom}(L)$}

C++ does not define structural subtyping\cite{stroustrup1994design} ((ISO/IEC 14882) ) in core type system: type compatibility is not determined by the structure but rather by explicit declarations, and while templates and concepts enable compile-time constraints resembling structural checks, they do not constitute formal structural subtyping judgments in the specification. 
C++ is fundamentally nominally typed: subtype relationships are explicitly declared via class inheritance (e.g., \texttt{class B : public A}), and the standard formally specifies these relationships through inheritance and access control rules. 

\[
    I_{nom}(C++) = 1, \quad I_{struct}(C++) = 0, \quad S_{nom}(C++) = \frac{I_{struct}(L) + I_{nom}(L)}{2}
    \Rightarrow S_{nom}(C++) = 0.5
\]

\textbf{2) Variance Formalization Score $S_{var}(L)$}

C++ supports parametric polymorphism via templates, but variance (covariance or contravariance) is neither explicitly declared nor syntactically annotated; instead, any variance behavior arises implicitly through pointer/reference conversions and template instantiation rules, without first-class variance constructs. 
The language enforces type safety through well-defined template instantiation and conversion rules, but does not provide a formally specified or proven variance system; moreover, certain constructs demonstrate that variance is only partially restricted rather than formally sound \cite{pierce2002types}.

\[
    V_{exp}(C++) = 0.5, \quad V_{sound}(C++) = 0.5, \quad S_{var}(L) = \frac{V_{exp}(L) + V_{sound}(L)}{2}
    \Rightarrow S_{var}(C++) = 0.5
\]

\textbf{3) Inheritance–Subtyping Separation Score $S_{rel}(L)$}

First, the language does not provide constructs for declaring subtyping independently of inheritance (e.g., no native interface or structural typing mechanism separate from class inheritance).

$I_{dec}(\text{C++}) = 0$

Second, subtyping rules—such as implicit conversion from derived to base class pointers or references—are directly derived from inheritance semantics, meaning there is no independent subtyping relation.

$I_{ind}(\text{C++}) = 0$

Third, the specification does not formally distinguish inheritance and subtyping as separate relations; instead, substitutability is an emergent property of inheritance rules and type conversion mechanisms, without an explicit formal separation.

$I_{form}(\text{C++}) = 0$

\[
    S_{rel}(\text{C++}) = \frac{0 + 0 + 0}{3} = 0
\]

\textbf{Final metric value:}
\[
    TSSI(C++) = 0.3 \cdot S_{nom}(L) + 0.3 \cdot S_{var}(L) + 0.4 \cdot S_{rel}(L) = 
\]

\subsubsection{Abstraction and Encapsulation Mechanisms}

\textbf{1) Visibility Granularity Index $V(L)$}

The language provides three semantically distinct access levels: \texttt{private}, \texttt{protected}, and \texttt{public}, each formally specified to constrain member accessibility in class hierarchies and across translation units.

\[
V(C++) = \frac{3}{4} = 0.75
\]


\textbf{2) Module Isolation Strength $M(L)$}

The language provides two primary boundary mechanisms: namespaces and modules. 
Namespaces serve as organizational scopes but do not enforce strict access control, as all members are visible once declared; hence, they are not counted as strict boundaries. 
In contrast, C++20 modules introduce explicit export control via \texttt{export} declarations, forming well-defined interface boundaries that restrict visibility across translation units.

\begin{itemize}
    \item $b_{strict}(\text{C++}) = 1$
    \item $b_{total}(\text{C++}) = 2$
\end{itemize}

\[
M(C++) = 0.5
\]


\textbf{3) Interface Abstraction Capacity $I(L)$}

The language does not provide a dedicated \textit{interface} keyword; instead, abstraction is achieved primarily through abstract classes and, more recently, through \texttt{concepts}, which define compile-time behavioral constraints. 
Regarding object-defining constructs, C++ includes \texttt{class}, \texttt{struct}, and \texttt{union}, all formally specified as distinct type constructors.

\begin{itemize}
    \item $n_{abstraction_constructs}(\text{C++}) = 2$ 
    \item $n_{object}(\text{C++}) = 3$
\end{itemize}

\[
    I(C++) = \frac{2}{3} = 0.67
\]


\textbf{4) Encapsulation Enforcement Coefficient $E(L)$}

C++ enforces multiple mechanisms that allow bypassing encapsulation constraints: pointer arithmetic, type casting (e.g., \texttt{reinterpret\_cast}), direct memory manipulation, and undefined behavior that can expose or modify object internals without restriction. Furthermore, the absence of runtime enforcement and the presence of low-level operations mean that encapsulation is not strictly preserved under all conditions.

\[
    E(C++) = 0
\]

\textbf{Final metric value:}
\[
    AE(C++) = 0.2 \cdot V(L) + 0.25 \cdot M(L) + 0.25 \cdot I(L) + 0.3 \cdot E(L)
\]

\subsubsection{Inheritance Semantics and Hierarchy Management}

\textbf{1) Hierarchy Simplicity Factor $H_s(L)$}

C++ class may inherit from multiple base classes using a comma-separated base-specifier list, which implies that the maximum number of direct superclasses $M(L)$ is unbounded in practice (i.e., $M(L) > 1$ and limited only by implementation constraints). Since there is no fixed upper bound in the specification, C++ effectively reaches the maximal considered value in comparative analysis.

\[
    M(\text{C++}) = M_{max}
    H_s(C++) = 1 - \frac{M(C++) - 1}{M_{max} - 1} = 0
\]

\textbf{2) Linearization Determinism Score $L_d(L)$}

The C++ standard does not define a global, formally specified method resolution order (MRO) comparable to linearization algorithms, instead, name lookup and overload resolution are governed by a combination of rules including scope-based lookup, unqualified and qualified name lookup, and class member access control. In multiple inheritance scenarios, ambiguities may arise (e.g., the diamond problem), and the programmer must explicitly disambiguate member access using scope resolution or employ virtual inheritance to alter object layout and lookup behavior. Since the standard specifies lookup procedures but does not define a single deterministic linearization sequence for all cases, the effective resolution strategy is only partially specified.

\[
    L_d(C++) = 0.5
\]

\textbf{3) Constraint Protection Coefficient $C_p(L)$}

$p_1 = 1$: The language provides \texttt{final} for classes and virtual functions, preventing further inheritance or overriding.

$p_2 = 0$: It does not support sealed or closed hierarchies with an explicitly fixed set of subclasses.

$p_3 = 0$: The \texttt{override} specifier exists but is not mandatory, hence no enforced override requirement.

$p_4 = 1$: Member functions are non-virtual by default unless explicitly declared \texttt{virtual}.

$p_5 = 0$: Although abstract classes exist via pure virtual functions, there is no strict language-level separation of interfaces as a distinct construct.

$p_6 = 0$: Compile-time inheritance restrictions are partially present (e.g., \texttt{final}), but since this is already accounted for in $p_1$ and no additional independent restriction system exists. 

\[
    C_p(\text{C++}) = \frac{1 + 0 + 0 + 1 + 0 + 0}{6} \approx 0.33
\]

\textbf{4) Fragile Base Mitigation Factor $F_b(L)$}

$s_1 = 0$: The language does not require mandatory override annotations, as \texttt{override} is optional.

$s_2 = 1$: It provides strict method visibility control via access specifiers (\texttt{private}, \texttt{protected}, \texttt{public}), which constrain exposure of overridable members.

$s_3 = 0$: There is no requirement for explicit superclass method invocation when overriding (calls to base class implementations via qualified names are optional).

$s_4 = 1$: C++ enforces variance restrictions in overriding, allowing only limited covariance (e.g., return types for pointers/references) and disallowing unsafe contravariance, which ensures type safety at compile time.

$s_5 = 1$: Member functions are non-virtual by default and require explicit \texttt{virtual} to enable overriding.

$s_6 = 1$: Finally, the compiler performs strict compile-time checks on overriding consistency (e.g., signature matching, virtual dispatch rules, and diagnostics with \texttt{override}). 

\[
    F_b(\text{C++}) = \frac{0 + 1 + 0 + 1 + 1 + 1}{6} \approx 0.67
\]

\textbf{Final metric value:}
\[
    ISRI(C++) = 0.2 \cdot H_s(L) + 0.25 \cdot L_d(L) + 0.3 \cdot C_p(L) + 0.25 \cdot F_b(L)
\]

\subsubsection{Composition and Alternatives to Inheritance}

\textbf{1) Class-Based Delegation Score $D_{cls}(L)$}

There is no dedicated keyword or built-in construct (such as \texttt{delegate}) for delegation in C++. Instead, delegation must be implemented manually using composition and forwarding functions (e.g., wrapper or adapter patterns). Therefore, delegation is only emulated. 
Furthermore,  C++ does not provide automatic method forwarding mechanisms, meaning developers must explicitly define forwarding functions.

\[
    K_{del}(C++) = 0.5, \quad S_{fwd}(C++) = 0
    \Rightarrow D_{cls}(C++) = 0.25
\]

\textbf{2) Trait Expressiveness Score $T_{exp}(L)$}

In cases of multiple inheritance, name conflicts between base classes result in ambiguity errors unless explicitly resolved by the programmer using qualified names; however, the language does not provide trait-style exclusion or aliasing mechanisms. 
C++ templates and mixin patterns do not support explicit declaration of required methods or properties within the mixin itself.

\[
    T_{exp}(C++) = 0
\]

\textbf{3) Interface Composition Capacity $I_{comp}(L)$}

C++ allows a class to inherit from an arbitrary number of base classes, including abstract classes (i.e., classes with pure virtual functions), with no language-imposed upper bound. 
Unlike traditional interface-only constructs, abstract classes in C++ may contain fully defined member functions alongside pure virtual functions, effectively supporting default method implementations within interface-like structures. 

\[
    I_{comp}(C++) = 1
\]

\textbf{Final metric value:}
\[
    CRFI(C++) = 0.3 \cdot D_{cls}(L) + 0.35 \cdot T_{exp}(L) + 0.35 \cdot I_{comp}(L)
\]

\subsubsection{Object Representation and Creation Model}

\textbf{1) Identity Semantics Score $S_{id}(L)$}

The primary user-definable type categories include classes, structs, unions and enums. In C++, all these types follow value semantics by default: objects are copied on assignment and passed by value unless explicitly handled via pointers or references, which are not independent user-defined type categories but rather type modifiers.

\[
    S_{id}(C++) = \frac{0}{4} = 0
\]

\textbf{2) Instantiation Paradigm Diversity $S_{cr}(L)$}

First, C++ supports direct construction via constructors, including stack allocation and dynamic allocation using the \texttt{new} operator, which together represent a single class-based instantiation paradigm. Second, the language supports copy and move construction (including cloning via copy constructors), which forms a prototype-like paradigm grounded in object copying semantics. Third, C++ provides aggregate and uniform initialization (including brace-initialization), which constitutes a literal-style construction mechanism for certain types.

\[
    S_{cr}(C++) = \frac{3}{P_{max}}
\]

\textbf{3) Meta-level Extensibility Score $S_{meta}(L)$}

The language provides limited runtime introspection via mechanisms such as \texttt{typeid} and \texttt{dynamic\_cast}, which allow querying certain aspects of an object's dynamic type within polymorphic hierarchies. However, the language does not support structural modification of objects at runtime, nor does it provide a metaobject protocol for overriding object instantiation or altering class definitions dynamically. All type structures are fixed at compile time, and templates operate purely at compile time without enabling runtime intercession.

\[
    I_{level}(C++) = 1 \Rightarrow S_{meta}(C++) = \frac{1}{I_{max}}
\]

\textbf{Final metric value:}
\[
    ORCI(C++) = 0.35 \cdot S_{id}(L) + 0.3 \cdot S_{cr}(L) + 0.35 \cdot S_{meta}(L)
\]

\subsubsection{Polymorphism and Dispatch Mechanisms}

\textbf{1) Dispatch Scope Mechanism $M_{scope}(L)$}

First, C++ supports single dispatch via virtual functions, where method selection is determined by the dynamic type of the receiver object.

$I_{single}=1$

Second, C++ does not natively support multiple dispatch (i.e., dispatch based on the runtime types of multiple arguments); instead, such behavior must be emulated using patterns like double dispatch or visitor.

$I_{multi}=0$

Third, predicate-based dispatch (selection based on arbitrary runtime conditions or guards integrated into the dispatch mechanism) is not part of the language semantics and must be implemented manually using conditional logic.

$I_{pred}=0$

\[
    M_{scope}(C++) = \frac{1}{5} = 0.2
\]

\textbf{2) Dispatch Safety Guarantee $S_{disp}(L)$}

The compiler checks function signatures and virtual override consistency, dynamic dispatch via virtual functions depends on runtime object types and does not ensure complete safety.

$I_{static}(\text{C++}) = 0$

Furthermore, C++ does not enforce exhaustive dispatch definitions: constructs such as virtual method hierarchies or \texttt{switch} statements over enums do not require covering all possible dynamic types or cases at compile time, and missing cases are not universally diagnosed as errors.

$I_{exhaust}(\text{C++}) = 0$

\[
    S_{disp}(\text{C++}) = \frac{0 + 0}{2} = 0
\]

\textbf{3) Virtualization Optimization Potential $O_{virt}(L)$}

First, C++ supports non-virtual methods by default and provides the \texttt{final} specifier for both classes and virtual functions, preventing further overriding and enabling compiler optimizations.

$K_{final}=1$

Second, the \texttt{final} specifier applied to classes effectively creates sealed hierarchies by prohibiting inheritance.

$K_{sealed}=1$

Third, C++ supports static dispatch through non-virtual member functions, free functions, and templates, all of which are resolved at compile time and provide strong hints for devirtualization.

$K_{static}=1$

\[
    O_{virt}(C++) = \frac{1 + 1 + 1}{3} = 1
\]

\textbf{Final metric value:}
\[
    PDEI(C++) = 0.35 \cdot M_{scope}(L) + 0.35 \cdot S_{disp}(L) + 0.3 \cdot O_{virt}(L)
\]

\subsubsection{State Model, Mutability, and Concurrency Guarantees}

\textbf{1) Default Mutability Score $M_{def}(L)$}

By default, any variable declared without the \texttt{const} qualifier permits modification of its stored value after initialization, which directly implies that mutability does not require any additional annotation, whereas immutability does. The \texttt{const} qualifier must be explicitly applied to enforce immutability, and this requirement is purely opt-in rather than default behavior.

\[
    M_{def}(C++) = 0
\]

\textbf{2) Static Reference Constraint Density $S_{dens}(L)$}

In C++, the set of reference-related constructs includes syntactic forms that introduce or manipulate references and pointers as defined by the ISO C++ standard: namely raw pointers (\texttt{T\*}), lvalue references (\texttt{T\&}), rvalue references (\texttt{T\&\&}), pointer qualifiers such as \texttt{const} and \texttt{volatile}, smart pointers (e.g., \texttt{std\:\:unique\_ptr}, \texttt{std\:\:shared\_ptr}), and reference-related type modifiers in templates and type deduction.

$N_{ref}(\text{C++}) = 6$

Among these, constructs that impose static constraints include: (1) \texttt{const} qualification, which enforces compile-time immutability constraints on referenced objects; (2) lvalue references, which must bind to valid lvalues and cannot be reseated; (3) rvalue references, which restrict binding to temporaries and enable move semantics; and (4) \texttt{std::unique\_ptr}, which enforces exclusive ownership semantics at compile time via deleted copy constructors.

$N_{constr}(\text{C++}) = 4$

\[
    S_{dens}(C++) \approx 0.67
\]

\textbf{3) Concurrency Primitive Safety $P_{conc}(L)$}

In C++, concurrency semantics are defined by the ISO C++ standard (e.g., §1.10 and §30 in C++17 and later), where the memory model formally specifies data races as undefined behavior but does not provide compile-time guarantees of race freedom. The language relies on programmer discipline, supported by low-level primitives such as \texttt{std::thread}, \texttt{std::mutex}, and \texttt{std::atomic}, without a type system capable of proving absence of data races.

$S_{race}(\text{C++}) = 0$

Regarding abstraction level, C++ primarily offers low-level threading and synchronization constructs, although it also includes \texttt{std::future} and \texttt{std::async}, which provide limited structured concurrency; however, these do not constitute a fully structured or high-level concurrency model such as actors or CSP channels, and the dominant paradigm remains explicit threads and locks.

$S_{model}(C++) = 0$

\[
    P_{conc}(C++) = 0
\]

\textbf{Final metric value:}
\[
    SCSI(C++) = 0.35 \cdot M_{def}(L) + 0.25 \cdot S_{dens}(L) + 0.4 \cdot P_{conc}(L)
\]

\subsubsection{Support for Formal Specification and Verifiability}

\textbf{1) Native Contract Constructs Score $C_{nat}(L)$}

The ISO C++ standard  does not include fully adopted and semantically enforced language-level support for contracts. Consequently, there are no first-class syntactic constructs for specifying preconditions, postconditions, or invariants within the core language; instead, developers rely on mechanisms such as \texttt{assert}, comments, or external libraries, which do not satisfy the requirement of being both syntactically native and semantically defined by the language specification.

\[
    C_{nat}(C++) = 0
\]

\textbf{2) Verification Enforcement Level $V_{enf}(L)$}

In C++, verification enforcement must be assessed based on mechanisms defined by the ISO C++ standard, which does not provide built-in facilities for formal static verification. While the language supports runtime checking mechanisms such as \texttt{assert} and exception handling, these operate exclusively at runtime and are not part of a formal verification system guaranteed by the compiler.

\[
    V_{enf}(C++) = 0.5
\]

\textbf{Final metric value:}
\[
    FSVI(C++) = 0.45 \cdot C_{nat}(L) + 0.55 \cdot V_{enf}(L)
\]

\subsection{Metrics Evaluation for C}

\subsubsection{Introduction}

The C programming language, standardized by ISO/IEC 9899, represents a minimalistic procedural paradigm with a weak notion of abstraction and no native support for object-oriented constructs. As noted by Kernighan and Ritchie, “C provides no direct support for object-oriented programming” \cite{ritchie1978c}. Consequently, evaluating C under object-model-oriented metrics requires interpreting its type system, memory model, and modularization mechanisms through the lens of emulated object-oriented patterns. This subsection computes the defined metrics strictly based on the formal language specification and authoritative analyses \cite{iso9899_2024, pierce2002types}.

\subsubsection{1) Type System and Subtyping Model (TSSI)}

\textbf{1) Nominal–Structural Typing Score $S_{nom}(L)$}

Structural subtyping is absent in C, since type compatibility is not based on structural equivalence with subtyping rules, but rather on strict compatibility rules without a formal subtype relation (ISO/IEC 9899).

$I_{struct}(C)=0$. 

Nominal typing with explicit subtype declarations is also absent: although C provides named types via \texttt{struct}, \texttt{union}, and \texttt{typedef}, it does not support inheritance or explicit subtype relationships between named types, as required by nominal subtyping systems.

$I_{nom}(C)=0$.

\[
S_{nom}(C) = 0
\]

\textbf{2) Variance Formalization Score $S_{var}(L)$}

C does not support generics in the sense required for variance (i.e., no parameterized types with subtype relations), as its type system lacks constructs such as type parameters; therefore, variance annotations are inapplicable.

$V_{exp}(C)=0$. 
Since variance is not defined at all in the language specification, there are no formal rules governing covariance, contravariance, or invariance, nor any associated soundness proofs.

$V_{sound}(C)=0$.

\[
    S_{var}(C) = 0
\]

\textbf{3) Inheritance–Subtyping Separation Score $S_{rel}(L)$}

Since C defines neither inheritance nor subtyping:
\[
    E_{inh}(C) = \varnothing, \quad E_{sub}(C) = \varnothing
\]

C is a procedural language with a type system based on compatibility rather than subtype polymorphism; constructs such as \texttt{struct} composition and pointer conversions do not constitute explicit subtyping relations defined by the language semantics.

$I_{dec}(\mathrm{C}) = 0$ because the language does not provide explicit subtyping independent of inheritance (nor inheritance itself).

$I_{ind}(\mathrm{C}) = 0$ since no independent subtyping rules are specified.

$I_{form}(\mathrm{C}) = 0$ because there is no formal distinction between inheritance and subtyping in the specification. 

\[
    S_{rel}(\mathrm{C}) = \frac{0 + 0 + 0}{3} = 0
\]

\textbf{Final metric value:}
\[
    TSSI(C) = 0.3 \cdot S_{nom}(L) + 0.3 \cdot S_{var}(L) + 0.4 \cdot S_{rel}(L) = 
\]

\subsubsection{Abstraction and Encapsulation Mechanisms}

\textbf{1) Visibility Granularity Index $V(L)$}

C supports two semantically distinct visibility scopes via storage-class specifiers and linkage rules: \textit{external linkage} (identifiers visible across translation units) and \textit{internal linkage} (restricted to a single translation unit using the \texttt{static} keyword).

\[
    n_{levels}(C) = 2
\]

\[
    V(C) = \frac{2}{n_{max}}
\]


\textbf{2) Module Isolation Strength $M(L)$}

C provides a single primary modularization construct: the \textit{translation unit}, formed after preprocessing, which serves as the compilation and visibility boundary.
However, this boundary does not enforce explicit export control lists: visibility across translation units is governed implicitly by linkage rules (external vs. internal) and declarations in header files, rather than by a formal mechanism requiring explicit export specifications. Consequently, there are no boundaries with enforced export control.

\[
b_{strict}(C) = 0, \quad b_{total}(C) = 1
\Rightarrow M(C) = 0
\]


\textbf{3) Interface Abstraction Capacity $I(L)$}

C does not provide dedicated language-level constructs for behavioral abstraction such as interfaces, abstract classes, or protocols. Abstraction is achieved idiomatically using header files and function declarations, which are not distinct semantic constructs in the type system.

\[
    I(C) = 0
\]


\textbf{4) Encapsulation Enforcement Coefficient $E(L)$}

C does not enforce encapsulation at the language level: it allows unrestricted memory access via pointers, including arbitrary casting between incompatible types (subject only to undefined behavior constraints), direct field access to \texttt{struct} members, and manipulation of object representations through pointer arithmetic. Additionally, there are no runtime checks preventing access violations, nor controlled unsafe regions distinct from safe code.

\[
    E(C) = 0
\]

\textbf{Final metric value:}
\[
    AE(C) = 0.2 \cdot V(L) + 0.25 \cdot M(L) + 0.25 \cdot I(L) + 0.3 \cdot E(L)
\]

\subsubsection{Inheritance Semantics and Hierarchy Management}

According to the ISO C standard, C does not support inheritance:

\[
    M(C) = 0 \Rightarrow H_s(C) = 0
    L_d(C) = 0
    C_p(C) = 0
    F_b(C) = 0
\]

\textbf{Final metric value:}
\[
    ISRI(C) = 0.2 \cdot H_s(L) + 0.25 \cdot L_d(L) + 0.3 \cdot C_p(L) + 0.25 \cdot F_b(L)
\]

\subsubsection{Composition and Alternatives to Inheritance}

\textbf{1) Class-Based Delegation Score $D_{cls}(L)$}

C provides neither keywords nor built-in constructs for delegation. However, delegation-like behavior can be manually emulated using function pointers and struct composition (e.g., forwarding calls through embedded structures), which corresponds to adapter-based emulation. 
The language requires explicit function calls and does not include any facility for automatic method delegation or forwarding without boilerplate code. 

\[
    K_{del}(C) = 0.5, \quad S_{fwd}(C) = 0
    \Rightarrow D_{cls}(C) = 0.25
\]

\textbf{2) Trait Expressiveness Score $T_{exp}(L)$}

Since C does not support traits or multiple composition units with overlapping method namespaces, such conflicts cannot be resolved via language-defined mechanisms and instead typically result in compilation errors (e.g., duplicate symbol definitions). 
C lacks traits or mixins entirely, and thus provides no facility to specify required methods or fields for composition. 

\[
    T_{exp}(C) = 0
\]

\textbf{3) Interface Composition Capacity $I_{comp}(L)$}

No interfaces:
\[
I_{comp}(C) = 0
\]

\textbf{Final metric value:}
\[
    CRFI(C) = 0.3 \cdot D_{cls}(L) + 0.35 \cdot T_{exp}(L) + 0.35 \cdot I_{comp}(L)
\]

\subsubsection{Object Representation and Creation Model}

\textbf{1) Identity Semantics Score $S_{id}(L)$}

The set of user-definable type categories includes \texttt{struct}, \texttt{union}, and \texttt{enum}. These types are defined with value semantics: objects of these types are copied on assignment and passed by value unless explicitly manipulated via pointers, meaning the language does not enforce reference semantics for any user-defined type category.

\[
T_{user}(C) = 3, \quad T_{ref}(C) = 0
\Rightarrow S_{id}(C) = 0
\]

\textbf{2) Instantiation Paradigm Diversity $S_{cr}(L)$}

C does not provide high-level instantiation constructs such as \texttt{new}, prototype cloning, or actor-based spawning; instead, it supports two fundamental paradigms: (1) static or automatic allocation (e.g., stack or global variables) and (2) dynamic allocation via standard library functions such as \texttt{malloc}. These mechanisms correspond to low-level memory allocation rather than abstract object instantiation models, but they are the only native paradigms available in the language.

\[
    S_{cr}(C) = \frac{2}{P_{max}}
\]

\textbf{3) Meta-level Extensibility Score $S_{meta}(L)$}

C does not provide any built-in mechanisms for runtime introspection or structural modification of objects; all type information is resolved at compile time, and no standardized runtime type system exists.

\[
    I_{level}(C) = 0 \Rightarrow S_{meta}(C) = 0
\]

\textbf{Final metric value:}
\[
    ORCI(C) = 0.35 \cdot S_{id}(L) + 0.3 \cdot S_{cr}(L) + 0.35 \cdot S_{meta}(L)
\]

\subsubsection{Polymorphism and Dispatch Mechanisms}

\textbf{1) Dispatch Scope Mechanism $M_{scope}(L)$}

C does not support single dispatch in the object-oriented sense, as there are no methods or implicit receiver-based calls; function calls are resolved statically at compile time.

$I_{single} = 0$

Similarly, C does not provide multiple dispatch (i.e., dynamic selection based on multiple argument types).

$I_{multi} = 0$

Predicate-based dispatch is also not natively supported; although conditional logic (e.g., \texttt{if}, \texttt{switch}) can be used to manually implement behavior selection, this is not part of the language's dispatch semantics.

$I_{pred} = 0$

\[
    M_{scope}(C) = 0
\]

\textbf{2) Dispatch Safety Guarantee $S_{disp}(L)$}

C does not provide built-in polymorphic dispatch mechanisms. Instead, dispatch is performed either statically via direct function calls or manually through function pointers. Although direct calls are type-checked, the use of function pointers does not guarantee full compile-time verification of dispatch correctness.

$I_{static}(\mathrm{C}) = 0$

C does not enforce exhaustive dispatch definitions: constructs like \texttt{switch} statements do not require covering all possible cases (even for enumerations), and missing cases are not mandated to be diagnosed by the compiler.

$I_{exhaust}(\mathrm{C}) = 0$

\[
    S_{disp}(\mathrm{C}) = \frac{0 + 0}{2} = 0
\]

\textbf{3) Virtualization Optimization Potential $O_{virt}(L)$}

Since C does not define virtual methods or dynamic dispatch, the notion of \texttt{final} methods is not applicable.

$K_{final} = 0$

Likewise, C lacks class hierarchies entirely, and therefore provides no mechanism for restricting subclassing (i.e., sealed hierarchies).

$K_{sealed} = 0$

However, all function calls in C are statically bound by default (except via function pointers), and the \texttt{static} keyword allows internal linkage and potential compiler optimizations, which can be interpreted as a form of static dispatch hint within the language specification.

$K_{static} = 1$ 

\[
    O_{virt}(C) = \frac{1}{3}
\]

\textbf{Final metric value:}
\[
    PDEI(C) = 0.35 \cdot M_{scope}(L) + 0.35 \cdot S_{disp}(L) + 0.3 \cdot O_{virt}(L)
\]

\subsubsection{State Model, Mutability, and Concurrency Guarantees}

\textbf{1) Default Mutability Score $M_{def}(L)$}

In C, all objects are mutable by default unless explicitly qualified with the \texttt{const} type qualifier, as defined in the ISO C standard (e.g., ISO/IEC 9899), which states that the absence of \texttt{const} permits modification of stored values through lvalues.

\[
    M_{def}(C) = 0
\]

\textbf{2) Static Reference Constraint Density $S_{dens}(L)$}

In C, references are represented exclusively via pointers: (1) pointer types (\texttt{T*}), (2) pointer qualifiers such as \texttt{const}, \texttt{volatile}, and \texttt{restrict}, and (3) array-to-pointer decay in expressions.

$N_{ref}(\mathrm{C}) = 3$ 

Next, distinct syntactic constructs include: the \texttt{const} qualifier enforces compile-time immutability of the referenced object, and \texttt{restrict} imposes aliasing constraints that enable compiler optimizations by guaranteeing exclusive access within a scope; \texttt{volatile} does not impose aliasing or mutability constraints but rather affects optimization semantics, and pointer types and decay introduce no restrictions.

$N_{constr}(\mathrm{C}) = 2$

\[
    S_{dens}(C) = \frac{2}{3}
\]

\textbf{3) Concurrency Primitive Safety $P_{conc}(L)$}

First, C provides no type-system or runtime enforcement of data race freedom. Although the standard defines a memory model and introduces atomic operations via \texttt{\_Atomic} and \texttt{stdatomic.h}, correct synchronization is entirely the programmer’s responsibility, and violations result in undefined behavior.

$S_{race}(\mathrm{C}) = 0$

Second, C only offers low-level concurrency primitives, such as threads (via \texttt{threads.h}) and mutexes, without higher-level abstractions like actors, channels, or structured task systems.

$S_{model}(\mathrm{C}) = 0$

\[
    P_{conc}(C) = 0
\]

\textbf{Final metric value:}
\[
    SCSI(C) = 0.35 \cdot M_{def}(L) + 0.25 \cdot S_{dens}(L) + 0.4 \cdot P_{conc}(L)
\]

\subsubsection{Support for Formal Specification and Verifiability}

\textbf{1) Native Contract Constructs Score $C_{nat}(L)$}

The language provides no native support for design-by-contract features: preconditions and postconditions are not part of the function signature or type system, and invariants are not expressible as built-in semantic constructs.

\[
    C_{nat}(C) = 0
\]

\textbf{2) Verification Enforcement Level $V_{enf}(L)$}

C does not support static verification in the sense of compile-time proof of program correctness; its type system is not expressive enough to encode and verify formal specifications.

\[
    V_{enf}(C) = 0
\]

\textbf{Final metric value:}
\[
    FSVI(C) = 0.45 \cdot C_{nat}(L) + 0.55 \cdot V_{enf}(L)
\]
